\subsection{Описание метода}

\subsubsection{Постановка задачи Коши для ОДУ}

\textbf{Обыкновенным дифференциальным уравнением} называется уравнение, которое связывает независимую переменную $x$, искомую функцию $y(x)$ и её производные различных порядков.

Если в уравнение входит только первая производная, оно называется ОДУ \textit{первого порядка}:
\[
y' = f(x, y).
\]

\textbf{Задачей Коши} \textit{(или задачей с начальными условиями)} для дифференциального уравнения первого порядка называется задача поиска функции $y(x)$, которая удовлетворяет самому уравнению на интервале $[a, b]$, а также заданному \textit{начальному условию}:
\[
y(a) = y_0,
\]
где $y_0$ --- известное начальное значение в точке $x = a$. 

\subsubsection{Метод Эйлера}

Рассмотрим простейший одношаговый численный метод --- \textbf{метод Эйлера}. Пусть на отрезке $[a, b]$ задан некоторый постоянный шаг $h = (b-a)\slash n$, тогда с его помощью определим узлы сетки:
\[
  x_i=a+ih,\quad i=\overline{1,n}  
\]

Разложим истинное решение $y(x)$ в окрестности точки $x_i$ в ряд Тейлора:
\[
y(x_{i+1}) = y(x_i + h) = y(x_i) + h y'(x_i) + \frac{h^2}{2} y''(\xi_i), \quad \xi_i \in (x_i, x_{i+1}).
\]

Метод Эйлера строит строит точки \( \omega_i\approx y(x_i) \), отбрасывая члены высших порядков малости. Учитывая, что $y'(x_i) = f(x_i, y(x_i))$, имеем:
\[
  \begin{aligned}
    \omega_0&=y_0 \\
    \omega_{i+1} &= \omega_i + h f(x_i, \omega_i), \quad i = \overline{0,n-1}
  \end{aligned}
\]

Для поиска \textit{локальной погрешности} метода предположим, что в узле $x_i$ вычисленное значение совпадает с точным, то есть $\omega_i = y(x_i)$. Разложим истинное решение $y(x_{i+1})$ в окрестности точки $x_i$ в ряд Тейлора:
\[
y(x_{i+1}) = y(x_i) + h \cdot y'(x_i) + \frac{h^2}{2} \cdot y''(\xi_i), \quad \xi_i \in (x_i, x_{i+1}).
\]
По определению ОДУ и методу Эйлера:
\[
  y'(x_i) = f(x_i, y(x_i))\implies |y(x_{i+1}) - \omega_{i+1}| = \frac{h^2}{2} |y''(\xi_i)| = O(h^2)\ \square
\]

Чтобый найти \textit{глобальную погрешность}, учтём, что при работе на всем отрезке $[a, b]$ ошибки на каждом шаге суммируются. Чтобы пройти весь интервал, алгоритму требуется совершить $n$ шагов.

В худшем случае, если локальные погрешности накапливаются с одинаковым знаком, полная накопленная ошибка пропорциональна произведению числа шагов на локальную ошибку одного шага:
\[
|y(b) - y_n| \approx n \cdot O(h^2) = \frac{b - a}{h} \cdot O(h^2) = (b - a) \cdot O(h) = O(h)\ \square
\]

Таким образом, метод Эйлера обладает \textbf{первым порядком точности}.

\subsubsection{Критерий останова}

Чтобы достичь погрешности решения не более чем \( \varepsilon \), воспользуемся следующим наблюдением\cite{wikipedia_runge_rule}. Метод Эйлера имеет первый порядок точности. По определению это означает, что приближенное решение $y_h(b)$, полученное с шагом $h$, связано с точным решением $y(b)$ соотношением:
\[
y(b) = y_h(b) + C \cdot h + O(h^2),
\]
где $C$ --- константа, которая не зависит от $h$.

Если мы уменьшим шаг интегрирования в два раза и найдем решение $y_{h/2}(b)$, то для него справедливо:
\[
y(b) = y_{h/2}(b) + C \cdot \frac{h}{2} + O(h^2)
\]

Вычтем из второго равенства первое, получив следующее соотношение:
\[
\abs{y_{h/2}(b) - y_h(b)} \approx \abs{C \cdot \frac{h}{2}}
\]

Нехитрой подстановкой во второе равенство третьего заключим:
\[
  \abs{y(b) - y_{h/2}(b)} \approx \abs{y_{h/2}(b) - y_h(b)}
\]

Таким образом, для достижения точности $\varepsilon$ алгоритм должен итерационно уменьшать шаг $h$ до тех пор, пока не выполнится критерий останова:
\[
|y_{h/2}(b) - y_h(b)| \le \varepsilon
\]